\label{chap:p3-83-47-godel_cohen}

\section{De la cohérence constructible à l'indépendance par forcing}

\begin{theorem}[Gödel, 1940 : moitié constructible]
Dans la métathéorie usuelle,
\[
  \operatorname{Con}(\mathsf{ZF})
  \Longrightarrow
  \operatorname{Con}
  \bigl(\mathsf{ZF}+\mathsf{AC}+\mathsf{GCH}\bigr).
\]
Plus précisément, étant donné un modèle \(M\models\mathsf{ZF}\), son univers
constructible interne \(L^M\) satisfait, au sens interne correspondant,
\[
  \mathsf{ZF}+V=L+\mathsf{AC}+\mathsf{GCH}.
\]
\end{theorem}

\begin{proof}[Idée de la preuve classique]
La hiérarchie constructible est définie par récurrence sur les ordinaux,
en n'admettant à chaque étape que les sous-ensembles définissables sur les étapes
précédentes. La construction conserve les ordinaux, produit un modèle interne
et fournit un bon ordre définissable global, dont se déduisent le choix et
l'hypothèse généralisée du continu \citep{goedel1940}.
\end{proof}

\begin{theorem}[Cohen, 1963--1964 : moitié par forcing]
On a les implications métamathématiques
\[
  \operatorname{Con}(\mathsf{ZFC})
  \Longrightarrow
  \operatorname{Con}(\mathsf{ZFC}+\neg\mathsf{CH}),
\]
et
\[
  \operatorname{Con}(\mathsf{ZF})
  \Longrightarrow
  \operatorname{Con}(\mathsf{ZF}+\neg\mathsf{AC}).
\]
En combinant les deux moitiés, \(\mathsf{CH}\) est indépendante de
\(\mathsf{ZFC}\) et \(\mathsf{AC}\) est indépendant de \(\mathsf{ZF}\),
toujours sous l'hypothèse de cohérence correspondante
\citep{cohen1963,cohen1964}.
\end{theorem}

Pour la première construction, on part d'un modèle de base \(M\), d'un poset
\(\mathbb P\in M\) et d'un filtre \(M\)-générique \(G\subseteq\mathbb P\) :
\[
  G\cap D\ne\varnothing
  \qquad
  \text{pour tout dense }D\subseteq\mathbb P\text{ appartenant à }M.
\]
Le théorème du forcing relie \(p\Vdash_M\varphi\) à la vérité de
\(\varphi\) dans \(M[G]\). Pour obtenir
\(\mathsf{ZF}+\neg\mathsf{AC}\), on prend en outre un sous-modèle symétrique
approprié d'une extension générique ; une extension ordinaire d'un modèle de
\(\mathsf{ZFC}\) conserve le choix \citep{cohen1966}.

\begin{remark}[Conditions de validité]
Ces théorèmes prouvent la cohérence relative et l'indépendance. Ils ne prouvent pas la
cohérence absolue de \(\mathsf{ZF}\) ou \(\mathsf{ZFC}\), ne décident pas
\(\mathsf{CH}\) dans un univers privilégié et ne transforment pas tout filtre
compatible en filtre générique. Ils ne sélectionnent pas non plus les branches HMT ni
n'identifient une sortie numérique.
\end{remark}

\subsection{\(2\) opérations non équivalentes : extension du modèle et enrichissement de l'objet}

Le forcing et HMT n'effectuent pas la même opération. Le premier compare un
univers \(M\) à une extension \(M[G]\) ; la seconde conserve des fibres qu'un
lecteur oublie. Si \(\mathsf{Hist}_{\rm HMT}\) est la catégorie des histoires
enrichies, le foncteur
\[
  U:\mathsf{Hist}_{\rm HMT}\longrightarrow\mathsf{Set},
  \qquad
  (x_m,\mu_m)_{m\ge0}\longmapsto(x_m)_{m\ge0},
\]
peut identifier des histoires distinctes sans changer l'ensemble visible. La
question de Cohen est de savoir quels sous-ensembles existent dans le modèle ; la question
généalogique est de savoir combien de relèvements avec mémoire et opérateurs possède un objet
visible. Ce sont des axes compatibles, mais aucun ne remplace l'autre.

\begin{center}
\begin{tabularx}{.94\textwidth}{>{\bfseries}p{.22\textwidth}XX}
\toprule
Opération & Ce qui change & Ce qu'elle ne prouve pas à elle seule\\
\midrule
\(M\mapsto M[G]\) & Univers d'ensembles, de noms et de vérité forcée &
Unicité ou généalogie d'une section HMT.\\
\(\xi\mapsto U(\xi)\) & Quantité de mémoire visible dans la projection &
\(\mathsf{CH}\), \(\mathsf{AC}\) ou généricité sur \(M\).\\
\bottomrule
\end{tabularx}
\end{center}

\subsection{L'ensemble des parties HMT ne présuppose pas \texorpdfstring{\(\mathsf{CH}\)}{CH}}

Dans chaque modèle où \(X\) est un espace de Cantor,
\[
  |\Clop(X)|=\aleph_0,\qquad
  |\Closed(X)|=\mathfrak c,\qquad
  |\mathcal P(X)|=2^{\mathfrak c},
  \qquad \mathfrak c=2^{\aleph_0}.
\]
\(\mathsf{CH}\) affirme uniquement \(\mathfrak c=\aleph_1\). Par le théorème
de Cantor, \(\mathfrak c<2^{\mathfrak c}\) dans chacun des modèles
comparés. Lors du passage de \(M\) à \(M[G]\), de nouvelles branches, de nouveaux
codes fermés et de nouveaux sous-ensembles peuvent apparaître, bien que les cylindres finis
conservent leur règle de formation. L'ensemble complet des parties est sensible au modèle ; la
distinction ouvert-fermé--fermé--complet ne dépend pas de la décision de \(\mathsf{CH}\).

\section{Poset de préfixes compatibles et limite inverse des extensions}

On définit
\[
  \mathbb P_{\TPK}=\bigsqcup_{n\ge0}\mathcal S_n
\]
et nous ordonnons par prolongement : \(q\le p\) si \(q\) contient plus d'information
que \(p\). Pour chaque \(N\),
\[
  D_N=\{p\in\mathbb P_{\TPK}:\ell(p)\ge N\}.
\]
Si tout préfixe survivant peut être prolongé, chaque \(D_N\) est dense.

\begin{theorem}[Branches et filtres de profondeur]
Les sections de \(\OmegaInf\) sont en bijection avec les filtres compatibles
de \(\mathbb P_{\TPK}\) qui rencontrent tous les \(D_N\).
\end{theorem}

\begin{proof}
Les préfixes d'une section sont compatibles, dirigés et atteignent toute
profondeur. Réciproquement, un filtre qui rencontre chaque \(D_N\) détermine une
restriction unique à chaque niveau ; la propriété de direction exclut les conditions
incompatibles et produit une section.
\end{proof}

Rencontrer tous les denses \(D_N\) exprime une profondeur illimitée, et non
une généricité sur un modèle. Être \(M\)-générique exigerait de rencontrer tout
dense \(D\in M\).

\begin{proposition}[Cas de Cohen conditionnel]
Si \(\OmegaInf\) est compact, métrisable, de dimension zéro, parfait et sans
points isolés, il est homéomorphe à l'ensemble de Cantor ; le forcing des cylindres non vides
a une partie dense équivalente au forcing de Cohen
\citep{cohen1963,cohen1964}.
\end{proposition}

Une section unique produit, au contraire, un forcing atomique ou trivial.
Le mot « forcing » ne sélectionne pas à lui seul une branche.

Pour une section \(x=(x_m)_{m\ge0}\), ses préfixes forment un filtre \(G_x\)
qui traverse les denses de profondeur. Cependant,
\[
  \bigl(\forall N,\;G_x\cap D_N\ne\varnothing\bigr)
  \centernot\Longrightarrow
  \bigl(\forall D\in M\text{ dense},\;G_x\cap D\ne\varnothing\bigr).
\]
La première propriété exprime un prolongement illimité ; la seconde,
une généricité relative à un modèle. Une section compatible produite par
neuf phases du TPK ne devient pas, par nomenclature, un réel de Cohen.

\begin{remark}[Conditions de validité]
Réfutent une identification automatique avec Cohen : un cylindre terminal ou un
point isolé ; un filtre qui traverse tous les \(D_N\) mais omet un
dense de \(M\) ; deux mémoires compatibles sur la même branche visible ; ou
l'absence d'un plongement dense explicite. Le nom partagé « branche » n'est pas
un morphisme.
\end{remark}

\section{Correspondance entre sections généalogiques et classes de Borel de l'espace limite}

Les cylindres engendrent
\[
 \Gamma_{\rm HMT}^{\rm Bor}
 =
 \sigma\bigl(\{[p]_n:p\in\mathcal S_n\}\bigr).
\]
La classe est stable par complémentation, unions dénombrables et images réciproques
d'applications continues certifiées.

Si un jeu infini sur des histoires possède une condition de victoire borélienne,
le théorème de détermination borélienne de Martin garantit une stratégie
gagnante \citep{martin1975}. C'est une application classique à l'espace HMT ;
elle ne dérive ni AD, ni MM, ni une valeur de \(2^{\aleph_0}\), ni la catégoricité de
\(\mathcal P(\mathbb R)\).

Pour \(A,B\subseteq\OmegaInf\), la réduction de Wadge
\[
 A\le_WB
\]
exige une application continue \(f\) avec \(x\in A\iff f(x)\in B\). Si l'on
restreint les applications à des opérateurs TPK certifiés, il apparaît un préordre plus
étroit dont la bonne fondation ou la semi-linéarité nécessite une preuve propre.

\begin{remark}[Conditions de validité]
Une démonstration qui utilise la bonne fondation de Wadge pour prouver la
détermination qui justifierait ensuite cette bonne fondation est circulaire. Les
branches historiques AD/MM sont conservées comme architecture, et non comme théorèmes HMT.
\end{remark}

\section{Optimisation minimax sur des fenêtres finies de préfixes compatibles}

Soit \(P_N=\mathcal S_N\) l'ensemble fini des candidats, \(T_N\) un ensemble
fini de tests adversariaux et
\[
  A_N:P_N\times T_N\to[-1,1]
\]
un gain fixé avant le résultat.

\begin{theorem}[Minimax fini]
\[
 \max_{\sigma\in\Prob(P_N)}
 \min_{\tau\in\Prob(T_N)}
 \mathbb E_{\sigma,\tau}A_N
 =
 \min_{\tau\in\Prob(T_N)}
 \max_{\sigma\in\Prob(P_N)}
 \mathbb E_{\sigma,\tau}A_N.
\]
\end{theorem}

\begin{proof}
C'est la dualité forte entre le programme linéaire du sélecteur et celui de
l'adversaire, sous la forme du théorème minimax de von Neumann
\citep{vonneumann1928}.
\end{proof}

L'équilibre peut être mixte et ne fournit pas une branche pure unique.

\section{Système projectif de troncatures et construction de la section limite}

Soient
\[
 X=\varprojlim P_N,
 \qquad
 Y=\varprojlim T_N
\]
compacts, et soit \(A:X\times Y\to\mathbb R\) continu. Les espaces de
probabilités sont compacts convexes dans la topologie faible. Le théorème de
Sion \citep{sion1958} donne
\[
 \sup_{\mu\in\Prob(X)}\inf_{\nu\in\Prob(Y)}
 \iint A\,d\mu\,d\nu
 =
 \inf_{\nu\in\Prob(Y)}\sup_{\mu\in\Prob(X)}
 \iint A\,d\mu\,d\nu.
\]

\begin{corollary}
Un équilibre global ne sélectionne une branche pure que si la mesure du sélecteur
est une Dirac \(\delta_x\). L'existence d'un équilibre mixte n'implique pas
la concentration.
\end{corollary}

\begin{remark}[Conditions de validité]
Le passage à la limite échoue sans compacité, continuité ou cohérence
projective des mesures. Si le gain est défini en utilisant la cible qu'il doit
sélectionner, le minimax est alimenté par la cible et n'explique pas la route.
\end{remark}

\section{Extensions projectives et matricielles apportées par les entrelaceurs opératoriels}

Le forcing organise les prolongements possibles ; Borel classe les ensembles
d'histoires ; le minimax construit la robustesse face aux falsificateurs. Aucun ne remplace
la loi structurelle qui produit un survivant unique. Leur fonction est
de clarifier la topologie et la décision de l'espace des routes, et non d'introduire une
sélection rétrospective.

\begin{remark}[Séparation des résultats et portée]
Les théorèmes de Gödel et Cohen sont des \emph{théorèmes classiques appliqués}. Le
poset, la comparaison typée entre filtre de profondeur et filtre générique,
la classe de Borel HMT et le jeu sélecteur--falsificateur sont des
\emph{constructions explicites du présent développement}. Martin, von Neumann et Sion s'appliquent
ensuite, avec leurs hypothèses classiques. Les anciennes affirmations PSC/HTT sur
CH, AD et MM ne sont pas promues.
\end{remark}
